%% kapa: exam=1 round=33 year=2022 period=1 form=A subject=real_estate subject_ko=부동산학원론 seq=23 number=103
%% source: Q-Net 감정평가사 기출문제 (cst00309) artlSeq=5209997 file=raw/1008/5209997/2022년도 제33회 감정평가사 1차시험 1교시(민법,경제학,부동산학).pdf
%% conversion: kordoc 4.24.51 -> tools/gen1.py (auto); review=proofread-sonnet; flags=none
\begin{question}{103}
부동산관리와 생애주기에 관한 설명으로 옳지 않은 것은?
\begin{choices}
  \item 자산관리(Asset Management)란 소유자의 부를 극대화시키기 위하여 대상부동산을 포트폴리오 관점에서 관리하는 것을 말한다.
  \item 시설관리(Facility Management)란 각종 부동산시설을 운영하고 유지하는 것으로 시설 사용자나 건물주의 요구에 단순히 부응하는 정도의 소극적이고 기술적인 측면의 관리를 말한다.
  \item 생애주기상 노후단계는 물리적·기능적 상태가 급격히 악화되기 시작하는 단계로 리모델링을 통하여 가치를 올릴 수 있다.
  \item 재산관리(Property Management)란 부동산의 운영수익을 극대화하고 자산가치를 증진시키기 위한 임대차관리 등의 일상적인 건물운영 및 관리뿐만 아니라 부동산 투자의 위험관리와 프로젝트 파이낸싱 등의 업무를 하는 것을 말한다.
  \item 건물의 이용에 의한 마멸, 파손, 노후화, 우발적 사고 등으로 사용이 불가능할 때까지의 기간을 물리적 내용연수라고 한다.
\end{choices}
\end{question}
\answer{4}
